\documentclass[11pt,letterpaper]{article}
\usepackage{cornell-notes}
\usepackage{needspace}

\newcommand{\CornellDocumentTitle}{Chapter 11: unix and linux security}
\title{\CornellDocumentTitle}
\author{}
\date{}
\setDocTitle{\CornellDocumentTitle}
\setDocSubtitle{Cybersecurity Cornell Notes}
\setCornellCollection{Cybersecurity Reference Notes}
\setCornellUnitType{Chapter}
\setCornellUnitNumber{11}
\setCornellUnitTitle{unix and linux security}
\setCornellCompanionLabel{Cybersecurity study companion}

\begin{document}
\maketitle
\begin{CornellOverviewBox}{Chapter Orientation}
\textbf{Chapter orientation.} This chapter examines unix and linux security within the broader domain of overview of system and network security. Its central problem is how to reason about user, unix, access, and file while keeping security objectives, operating assumptions, evidence, and recovery connected. A practitioner should approach the material through protocol behavior, exposed services, trust boundaries, configuration, traffic visibility, and layered protection. Useful prerequisites include basic risk, threat, control, and systems concepts.
\end{CornellOverviewBox}
\section*{Learning Objectives}
\begin{itemize}
\item Explain the security purpose and scope of UNIX and Linux Security.
\item Identify the assets, actors, assumptions, and trust boundaries associated with user, unix, and access.
\item Distinguish Introduction from UNIX And Security and explain where their responsibilities intersect.
\item Analyze how failures in user, unix, and access can affect confidentiality, integrity, availability, privacy, or safety.
\item Evaluate relevant controls through protocol behavior, exposed services, trust boundaries, configuration, traffic visibility, and layered protection.
\item Audit an implementation using network diagrams, packet evidence, rulesets, service inventories, configuration baselines, alerts, and flow records.
\item Prioritize remediation according to exploitability, consequence, control strength, and recovery readiness.
\item Verify that corrective actions change observable risk rather than documentation alone.
\end{itemize}
\section*{Cornell Cue-and-Notes}
\Needspace{8\baselineskip}
\subsection*{Introduction}
\begin{longtable}{>{\RaggedRight\arraybackslash}p{0.29\textwidth}|>{\RaggedRight\arraybackslash}p{0.65\textwidth}}
\CornellHeader
\CornellNoteRow{What security problem does Introduction address?}{\textbf{Core idea.} Treat Introduction as a structured part of UNIX and Linux Security, with particular attention to unix, distribution, access, and user. \textbf{Analytical lens.} This topic establishes a component of the chapter's argument; connect its definitions and mechanisms to a testable security objective. For unix, distribution, and access, follow traffic and state across each boundary, including malformed, unauthenticated, and partially failed conditions. \textbf{Security reasoning.} Identify what must remain protected, who or what can alter the expected state, which assumptions cross a trust boundary, and how failure changes confidentiality, integrity, availability, privacy, or safety. \textbf{Application.} The practitioner should trace traffic, validate protocol state, minimize exposure, and test prevention and detection controls. \textbf{Evidence.} Seek network diagrams, packet evidence, rulesets, service inventories, configuration baselines, alerts, and flow records. Related subtopics include The Aims of System Security, Authentication, Authorization, and Availability.}
\end{longtable}
\begin{CornellSummaryBox}{Topic Summary}
Introduction is best remembered as the connection between unix, distribution, access, and user and the chapter's larger control objective. A defensible review states the assumption, expected behavior, evidence, failure consequence, and required response.
\end{CornellSummaryBox}
\Needspace{8\baselineskip}
\subsection*{UNIX And Security}
\begin{longtable}{>{\RaggedRight\arraybackslash}p{0.29\textwidth}|>{\RaggedRight\arraybackslash}p{0.65\textwidth}}
\CornellHeader
\CornellNoteRow{Which assumptions matter in UNIX And Security?}{\textbf{Core idea.} Treat UNIX And Security as a structured part of UNIX and Linux Security, with particular attention to access, unix, kernel, and control. \textbf{Analytical lens.} This topic establishes a component of the chapter's argument; connect its definitions and mechanisms to a testable security objective. For access, unix, and kernel, follow traffic and state across each boundary, including malformed, unauthenticated, and partially failed conditions. \textbf{Security reasoning.} Identify what must remain protected, who or what can alter the expected state, which assumptions cross a trust boundary, and how failure changes confidentiality, integrity, availability, privacy, or safety. \textbf{Application.} The practitioner should trace traffic, validate protocol state, minimize exposure, and test prevention and detection controls. \textbf{Evidence.} Seek network diagrams, packet evidence, rulesets, service inventories, configuration baselines, alerts, and flow records. Related subtopics include Integrity, Traditional UNIX Systems, Confidentiality, and Kernel Space Versus User Land.}
\end{longtable}
\begin{CornellSummaryBox}{Topic Summary}
UNIX And Security is best remembered as the connection between access, unix, kernel, and control and the chapter's larger control objective. A defensible review states the assumption, expected behavior, evidence, failure consequence, and required response.
\end{CornellSummaryBox}
\Needspace{8\baselineskip}
\subsection*{Basic UNIX Security Overview}
\begin{longtable}{>{\RaggedRight\arraybackslash}p{0.29\textwidth}|>{\RaggedRight\arraybackslash}p{0.65\textwidth}}
\CornellHeader
\CornellNoteRow{How should a reviewer evaluate Basic UNIX Security Overview?}{\textbf{Core idea.} Treat Basic UNIX Security Overview as a structured part of UNIX and Linux Security, with particular attention to user, file, kernel, and access. \textbf{Analytical lens.} This topic is architecture-centered: locate components, interfaces, trust boundaries, dependencies, control points, and failure propagation. For user, file, and kernel, follow traffic and state across each boundary, including malformed, unauthenticated, and partially failed conditions. \textbf{Security reasoning.} Identify what must remain protected, who or what can alter the expected state, which assumptions cross a trust boundary, and how failure changes confidentiality, integrity, availability, privacy, or safety. \textbf{Application.} The practitioner should trace traffic, validate protocol state, minimize exposure, and test prevention and detection controls. \textbf{Evidence.} Seek network diagrams, packet evidence, rulesets, service inventories, configuration baselines, alerts, and flow records. Related subtopics include Semantics of User Space Security, Standard File and Device Access Semantics, Sticky Bit, and Read, Write, Execute.}
\end{longtable}
\begin{CornellSummaryBox}{Topic Summary}
Basic UNIX Security Overview is best remembered as the connection between user, file, kernel, and access and the chapter's larger control objective. A defensible review states the assumption, expected behavior, evidence, failure consequence, and required response.
\end{CornellSummaryBox}
\Needspace{8\baselineskip}
\subsection*{Achieving UNIX Security}
\begin{longtable}{>{\RaggedRight\arraybackslash}p{0.29\textwidth}|>{\RaggedRight\arraybackslash}p{0.65\textwidth}}
\CornellHeader
\CornellNoteRow{Why can Achieving UNIX Security fail in practice?}{\textbf{Core idea.} Treat Achieving UNIX Security as a structured part of UNIX and Linux Security, with particular attention to access, user, file, and unix. \textbf{Analytical lens.} This topic establishes a component of the chapter's argument; connect its definitions and mechanisms to a testable security objective. For access, user, and file, follow traffic and state across each boundary, including malformed, unauthenticated, and partially failed conditions. \textbf{Security reasoning.} Identify what must remain protected, who or what can alter the expected state, which assumptions cross a trust boundary, and how failure changes confidentiality, integrity, availability, privacy, or safety. \textbf{Application.} The practitioner should trace traffic, validate protocol state, minimize exposure, and test prevention and detection controls. \textbf{Evidence.} Seek network diagrams, packet evidence, rulesets, service inventories, configuration baselines, alerts, and flow records. Related subtopics include Discretionary Versus Mandatory Access, Control, System Patching, and Establishing Secure Account Use.}
\end{longtable}
\begin{CornellSummaryBox}{Topic Summary}
Achieving UNIX Security is best remembered as the connection between access, user, file, and unix and the chapter's larger control objective. A defensible review states the assumption, expected behavior, evidence, failure consequence, and required response.
\end{CornellSummaryBox}
\Needspace{8\baselineskip}
\subsection*{Protecting User Accounts And Strengthening Authentication}
\begin{longtable}{>{\RaggedRight\arraybackslash}p{0.29\textwidth}|>{\RaggedRight\arraybackslash}p{0.65\textwidth}}
\CornellHeader
\CornellNoteRow{What security problem does Protecting User Accounts And Strengthening Authentication address?}{\textbf{Core idea.} Treat Protecting User Accounts And Strengthening Authentication as a structured part of UNIX and Linux Security, with particular attention to user, ssh, nis, and file. \textbf{Analytical lens.} This topic is control-centered: map each safeguard to a specific failure mode, then test independence, coverage, and bypass paths. For user, ssh, and nis, follow traffic and state across each boundary, including malformed, unauthenticated, and partially failed conditions. \textbf{Security reasoning.} Identify what must remain protected, who or what can alter the expected state, which assumptions cross a trust boundary, and how failure changes confidentiality, integrity, availability, privacy, or safety. \textbf{Application.} The practitioner should trace traffic, validate protocol state, minimize exposure, and test prevention and detection controls. \textbf{Evidence.} Seek network diagrams, packet evidence, rulesets, service inventories, configuration baselines, alerts, and flow records. Related subtopics include Controlling Account Access, Local Files, Network Information System, and Using Pluggable Authentication Modules to.}
\end{longtable}
\begin{CornellSummaryBox}{Topic Summary}
Protecting User Accounts And Strengthening Authentication is best remembered as the connection between user, ssh, nis, and file and the chapter's larger control objective. A defensible review states the assumption, expected behavior, evidence, failure consequence, and required response.
\end{CornellSummaryBox}
\Needspace{8\baselineskip}
\subsection*{Limiting Superuser Privileges}
\begin{longtable}{>{\RaggedRight\arraybackslash}p{0.29\textwidth}|>{\RaggedRight\arraybackslash}p{0.65\textwidth}}
\CornellHeader
\CornellNoteRow{Which assumptions matter in Limiting Superuser Privileges?}{\textbf{Core idea.} Treat Limiting Superuser Privileges as a structured part of UNIX and Linux Security, with particular attention to access, superuser, unix, and special. \textbf{Analytical lens.} This topic is identity-centered: distinguish identification, authentication, authorization, session state, privilege, and accountability. For access, superuser, and unix, follow traffic and state across each boundary, including malformed, unauthenticated, and partially failed conditions. \textbf{Security reasoning.} Identify what must remain protected, who or what can alter the expected state, which assumptions cross a trust boundary, and how failure changes confidentiality, integrity, availability, privacy, or safety. \textbf{Application.} The practitioner should trace traffic, validate protocol state, minimize exposure, and test prevention and detection controls. \textbf{Evidence.} Seek network diagrams, packet evidence, rulesets, service inventories, configuration baselines, alerts, and flow records. Related subtopics include Gaining Root Privileges With Su, Using Groups Instead of Root, and Using the sudo(1) Mechanism.}
\end{longtable}
\begin{CornellSummaryBox}{Topic Summary}
Limiting Superuser Privileges is best remembered as the connection between access, superuser, unix, and special and the chapter's larger control objective. A defensible review states the assumption, expected behavior, evidence, failure consequence, and required response.
\end{CornellSummaryBox}
\Needspace{8\baselineskip}
\subsection*{Securing Local And Network File Systems}
\begin{longtable}{>{\RaggedRight\arraybackslash}p{0.29\textwidth}|>{\RaggedRight\arraybackslash}p{0.65\textwidth}}
\CornellHeader
\CornellNoteRow{How should a reviewer evaluate Securing Local And Network File Systems?}{\textbf{Core idea.} Treat Securing Local And Network File Systems as a structured part of UNIX and Linux Security, with particular attention to file, protocols, user, and encrypted. \textbf{Analytical lens.} This topic establishes a component of the chapter's argument; connect its definitions and mechanisms to a testable security objective. For file, protocols, and user, follow traffic and state across each boundary, including malformed, unauthenticated, and partially failed conditions. \textbf{Security reasoning.} Identify what must remain protected, who or what can alter the expected state, which assumptions cross a trust boundary, and how failure changes confidentiality, integrity, availability, privacy, or safety. \textbf{Application.} The practitioner should trace traffic, validate protocol state, minimize exposure, and test prevention and detection controls. \textbf{Evidence.} Seek network diagrams, packet evidence, rulesets, service inventories, configuration baselines, alerts, and flow records. Related subtopics include Directory Structure and Partitioning for, Security, Employing Read-Only Partitions, and Finding Special Files.}
\end{longtable}
\begin{CornellSummaryBox}{Topic Summary}
Securing Local And Network File Systems is best remembered as the connection between file, protocols, user, and encrypted and the chapter's larger control objective. A defensible review states the assumption, expected behavior, evidence, failure consequence, and required response.
\end{CornellSummaryBox}
\Needspace{8\baselineskip}
\subsection*{Network Configuration}
\begin{longtable}{>{\RaggedRight\arraybackslash}p{0.29\textwidth}|>{\RaggedRight\arraybackslash}p{0.65\textwidth}}
\CornellHeader
\CornellNoteRow{Why can Network Configuration fail in practice?}{\textbf{Core idea.} Treat Network Configuration as a structured part of UNIX and Linux Security, with particular attention to access, unix, services, and interface. \textbf{Analytical lens.} This topic establishes a component of the chapter's argument; connect its definitions and mechanisms to a testable security objective. For access, unix, and services, follow traffic and state across each boundary, including malformed, unauthenticated, and partially failed conditions. \textbf{Security reasoning.} Identify what must remain protected, who or what can alter the expected state, which assumptions cross a trust boundary, and how failure changes confidentiality, integrity, availability, privacy, or safety. \textbf{Application.} The practitioner should trace traffic, validate protocol state, minimize exposure, and test prevention and detection controls. \textbf{Evidence.} Seek network diagrams, packet evidence, rulesets, service inventories, configuration baselines, alerts, and flow records. Related subtopics include Detecting and Disabling Standard UNIX, Services, Basic Network Setup, and Host-Based Firewall.}
\end{longtable}
\begin{CornellSummaryBox}{Topic Summary}
Network Configuration is best remembered as the connection between access, unix, services, and interface and the chapter's larger control objective. A defensible review states the assumption, expected behavior, evidence, failure consequence, and required response.
\end{CornellSummaryBox}
\Needspace{8\baselineskip}
\subsection*{Improving The Security Of Linux And UNIX Systems}
\begin{longtable}{>{\RaggedRight\arraybackslash}p{0.29\textwidth}|>{\RaggedRight\arraybackslash}p{0.65\textwidth}}
\CornellHeader
\CornellNoteRow{What security problem does Improving The Security Of Linux And UNIX Systems address?}{\textbf{Core idea.} Treat Improving The Security Of Linux And UNIX Systems as a structured part of UNIX and Linux Security, with particular attention to checklist, unix, yale, and firewall. \textbf{Analytical lens.} This topic establishes a component of the chapter's argument; connect its definitions and mechanisms to a testable security objective. For checklist, unix, and yale, follow traffic and state across each boundary, including malformed, unauthenticated, and partially failed conditions. \textbf{Security reasoning.} Identify what must remain protected, who or what can alter the expected state, which assumptions cross a trust boundary, and how failure changes confidentiality, integrity, availability, privacy, or safety. \textbf{Application.} The practitioner should trace traffic, validate protocol state, minimize exposure, and test prevention and detection controls. \textbf{Evidence.} Seek network diagrams, packet evidence, rulesets, service inventories, configuration baselines, alerts, and flow records.}
\end{longtable}
\begin{CornellSummaryBox}{Topic Summary}
Improving The Security Of Linux And UNIX Systems is best remembered as the connection between checklist, unix, yale, and firewall and the chapter's larger control objective. A defensible review states the assumption, expected behavior, evidence, failure consequence, and required response.
\end{CornellSummaryBox}
\Needspace{8\baselineskip}
\subsection*{Additional Resources}
\begin{longtable}{>{\RaggedRight\arraybackslash}p{0.29\textwidth}|>{\RaggedRight\arraybackslash}p{0.65\textwidth}}
\CornellHeader
\CornellNoteRow{Which assumptions matter in Additional Resources?}{\textbf{Core idea.} Treat Additional Resources as a structured part of UNIX and Linux Security, with particular attention to management, unix, webmin, and useful. \textbf{Analytical lens.} This topic establishes a component of the chapter's argument; connect its definitions and mechanisms to a testable security objective. For management, unix, and webmin, follow traffic and state across each boundary, including malformed, unauthenticated, and partially failed conditions. \textbf{Security reasoning.} Identify what must remain protected, who or what can alter the expected state, which assumptions cross a trust boundary, and how failure changes confidentiality, integrity, availability, privacy, or safety. \textbf{Application.} The practitioner should trace traffic, validate protocol state, minimize exposure, and test prevention and detection controls. \textbf{Evidence.} Seek network diagrams, packet evidence, rulesets, service inventories, configuration baselines, alerts, and flow records. Related subtopics include Nmap, Useful Tools, Webmin, and Local Configuration System.}
\end{longtable}
\begin{CornellSummaryBox}{Topic Summary}
Additional Resources is best remembered as the connection between management, unix, webmin, and useful and the chapter's larger control objective. A defensible review states the assumption, expected behavior, evidence, failure consequence, and required response.
\end{CornellSummaryBox}
\begin{CornellWarningBox}{Historical Context}
Some products, protocols, examples, threat observations, or legal references in this chapter are historically bounded. Retain the underlying threat model and control objective, but verify present applicability before treating a named implementation as current guidance.
\end{CornellWarningBox}
\begin{CornellOverviewBox}{Defensive Guidance}
Apply the chapter by making assets, authority, trust, expected behavior, and failure response explicit. In practice, trace traffic, validate protocol state, minimize exposure, and test prevention and detection controls. A control is not fully supported until its operation, monitoring, ownership, and recovery behavior are demonstrated with evidence.
\end{CornellOverviewBox}
\section*{Threat-and-Control Map}
\begingroup\scriptsize\setlength{\tabcolsep}{2pt}
\begin{longtable}{>{\RaggedRight\arraybackslash}p{0.135\textwidth}>{\RaggedRight\arraybackslash}p{0.145\textwidth}>{\RaggedRight\arraybackslash}p{0.155\textwidth}>{\RaggedRight\arraybackslash}p{0.145\textwidth}>{\RaggedRight\arraybackslash}p{0.16\textwidth}>{\RaggedRight\arraybackslash}p{0.17\textwidth}}
\toprule\rowcolor{Navy}\color{white}\textbf{Asset} & \color{white}\textbf{Threat / Failure} & \color{white}\textbf{Vulnerability} & \color{white}\textbf{Impact} & \color{white}\textbf{Detection / Evidence} & \color{white}\textbf{Control}\\\midrule
\endfirsthead
\toprule\rowcolor{Navy}\color{white}\textbf{Asset} & \color{white}\textbf{Threat / Failure} & \color{white}\textbf{Vulnerability} & \color{white}\textbf{Impact} & \color{white}\textbf{Detection / Evidence} & \color{white}\textbf{Control}\\\midrule
\endhead
Network services, communications, and connected hosts & Unauthorized access, interception, manipulation, or disruption & Exposed services, weak protocol assumptions, or unsafe configuration & Loss of confidentiality, integrity, availability, or control & Packet, flow, host, and alert telemetry correlated to expected behavior & Segmentation, hardened configuration, authentication, filtering, and monitoring\\\midrule
Assumptions surrounding user & Misuse or unexpected state transition & Trust in unix is not validated & A local weakness becomes a broader compromise path & Review evidence for access and test negative paths & Make trust explicit, constrain authority, and fail safely\\\midrule
Recovery capability and decision evidence & Control failure remains unnoticed or cannot be reversed & Monitoring, ownership, or restoration has not been exercised & Longer exposure and slower, less reliable recovery & Alert-to-response exercises and restoration tests & Assigned ownership, actionable telemetry, preserved evidence, and rehearsed recovery\\\midrule
\bottomrule
\end{longtable}
\endgroup
\section*{Practical Security Checklist}
\begin{itemize}[label=$\square$]
\item Define the in-scope assets, users, services, data, and operational dependencies for UNIX and Linux Security.
\item Document the expected behavior and security objective of user.
\item Map identities, privileges, trust boundaries, and externally controlled inputs associated with user and unix.
\item Record the threat or failure being addressed, its prerequisites, likely path, and credible consequences.
\item Inspect network diagrams, packet evidence, rulesets, service inventories, configuration baselines, alerts, and flow records.
\item Test both the intended path and negative paths, including invalid state, partial failure, excessive privilege, and unavailable dependencies.
\item Confirm preventive controls are complemented by actionable detection and a named response owner.
\item Exercise containment, restoration, and access; record actual recovery behavior and unresolved dependencies.
\item Validate exceptions, compensating controls, expiration dates, and accountable approval.
\item Preserve reproducible evidence and tie each finding to affected assets, risk, remediation, and verification criteria.
\end{itemize}
\begin{CornellSummaryBox}{Chapter Synthesis}
The chapter's principal lesson is that unix and linux security must be evaluated as an operating security system rather than an isolated product or statement. The most important failure patterns involve user, unix, access, and file, especially when ownership, boundaries, telemetry, or recovery remain implicit. A reviewer should connect architecture and behavior to network diagrams, packet evidence, rulesets, service inventories, configuration baselines, alerts, and flow records, prioritize findings by reachable consequence and control independence, and verify remediation through observable change. Within Overview of System and Network Security, this chapter contributes a reusable method: state the objective, expose assumptions, test failure, preserve evidence, and learn from operation.
\end{CornellSummaryBox}
\section*{Self-Test Questions}
\begin{enumerate}
\item What is the central security objective of UNIX and Linux Security?
\item How does Introduction contribute to the chapter's broader security argument?
\item Distinguish UNIX And Security from Basic UNIX Security Overview.
\item Which trust assumptions should be made explicit when evaluating user?
\item How could a weakness involving unix affect confidentiality, integrity, availability, privacy, or safety?
\item What evidence would demonstrate that controls surrounding access operate as intended?
\item Why should prevention, detection, response, and recovery be evaluated together in this domain?
\item Scenario: A reachable service accepts traffic across a boundary but its assumptions and monitoring are undocumented. What should the reviewer examine first, and why?
\item Scenario: A control associated with file passes a configuration check but fails during an operational exercise. How should the finding be interpreted and prioritized?
\item How should a practitioner distinguish enduring principles in this chapter from time-sensitive tools, examples, or implementation details?
\end{enumerate}
\Needspace{8\baselineskip}
\section*{Answer Key}
\begin{enumerate}
\item The objective is to protect the relevant assets by understanding protocol behavior, exposed services, trust boundaries, configuration, traffic visibility, and layered protection and by connecting controls to measurable risk reduction.
\item Introduction provides one part of the causal chain: it should identify expected behavior, dependencies, failure modes, and the evidence needed to justify confidence.
\item UNIX And Security and Basic UNIX Security Overview address different portions of the problem. A strong comparison states their scope, assumptions, enforcement points, evidence, and how a failure in one affects the other.
\item Reviewers should identify who controls user, what inputs or dependencies it trusts, where privilege changes, what happens during partial failure, and how those assumptions are tested.
\item A weakness in unix can propagate across trust boundaries. The actual impact depends on reachable assets, granted authority, control independence, visibility, and recovery capability.
\item Useful evidence includes network diagrams, packet evidence, rulesets, service inventories, configuration baselines, alerts, and flow records; configuration alone is insufficient when operation, monitoring, or recovery is part of the claim.
\item Prevention reduces probability, detection limits dwell time, response limits spread, and recovery restores objectives. Omitting one layer creates an unexamined failure mode.
\item Begin by establishing scope, ownership, expected behavior, and the violated assumption. Then trace traffic, validate protocol state, minimize exposure, and test prevention and detection controls, preserving evidence for prioritization and follow-up.
\item Treat the exercise as stronger evidence than a static check. Prioritize according to reachable impact, likelihood of real operating conditions, missing compensating controls, and recovery difficulty.
\item Retain the principle, threat model, and control objective; separately label technologies, legal references, product examples, and threat observations whose applicability depends on time or environment.
\end{enumerate}
\section*{Key-Term Recap}
\begin{longtable}{>{\RaggedRight\arraybackslash}p{0.27\textwidth}>{\RaggedRight\arraybackslash}p{0.66\textwidth}}
\toprule\rowcolor{Navy}\color{white}\textbf{Term} & \color{white}\textbf{Meaning}\\\midrule
\endfirsthead
\toprule\rowcolor{Navy}\color{white}\textbf{Term} & \color{white}\textbf{Meaning}\\\midrule
\endhead
\textbf{Authentication} & The process of establishing confidence that an asserted identity is genuine. \\\midrule
\textbf{Access Control} & Rules and enforcement mechanisms that restrict actions on resources to authorized subjects. \\\midrule
\textbf{Authorization} & The decision process that determines which authenticated subject may perform which action on a resource. \\\midrule
\textbf{Confidentiality} & The property that information is not disclosed to unauthorized parties. \\\midrule
\textbf{Integrity} & The property that information and systems remain accurate, complete, and protected from unauthorized modification. \\\midrule
\textbf{Attack Surface} & The collection of reachable interfaces, inputs, privileges, dependencies, and operational paths an adversary might exploit. \\\midrule
\textbf{Availability} & The property that authorized users can access information and services when required. \\\midrule
\textbf{Metadata} & Descriptive information about data or activity that can support operations, investigations, or privacy-invasive inference. \\\midrule
\textbf{Encryption} & A keyed transformation of plaintext into ciphertext to protect confidentiality. \\\midrule
\textbf{Privacy} & The appropriate governance of information about people, including collection, use, disclosure, retention, and individual interests. \\\midrule
\bottomrule
\end{longtable}
\begin{CornellExamBox}{One-Sentence Takeaway}
Treat unix and linux security as a verifiable chain from assets and assumptions through controls, evidence, response, and recovery.
\end{CornellExamBox}
\end{document}
